% Generated draft: requires mathematical and source review.
\documentclass[11pt]{article}
\usepackage[margin=25mm]{geometry}
\usepackage{fontspec}
\setmainfont{Latin Modern Roman}
\usepackage{amsmath,amssymb,mathtools,unicode-math}
\setmathfont{Latin Modern Math}
\usepackage{enumitem,needspace,xurl,hyperref}
\setlength{\parindent}{0pt}
\setlength{\parskip}{5pt}
\setlength{\emergencystretch}{3em}
\newcommand{\N}{\mathbb{N}}
\newcommand{\Z}{\mathbb{Z}}
\newcommand{\Q}{\mathbb{Q}}
\newcommand{\R}{\mathbb{R}}
\newcommand{\C}{\mathbb{C}}
\newcommand{\F}{\mathbb{F}}
\newcommand{\A}{\mathbb{A}}
\newcommand{\PP}{\mathbb{P}}
\newcommand{\E}{\mathbb{E}}
\newcommand{\eps}{\varepsilon}
\newcommand{\dd}{\,\mathrm{d}}
\begin{document}
\section*{UnsolvedMath ID: 30002140}
\textbf{Source problem code:} catalogue entry
\subsection{Definitions and mathematical statement}
Let \(\mathbb H^n\) denote real hyperbolic \(n\)-space of constant sectional curvature \(-1\), and let \(G_n=\operatorname{Isom}(\mathbb H^n)\) be its full isometry group. The source does not specify the permitted dimensions \(n\).

For \(\kappa\in\{0,-1\}\), a geodesic metric space \((X,d)\) is \(\operatorname{CAT}(\kappa)\) if every geodesic triangle satisfies the comparison inequality with its comparison triangle in the complete simply connected surface of constant curvature \(\kappa\): the distance between any two points on the triangle is at most the distance between the corresponding points on the comparison triangle. Corresponding points have the same distances from the respective endpoints along corresponding sides. For these nonpositive curvature parameters, no perimeter restriction is required.

An isometric action of \(G_n\) on \(X\) is a homomorphism \(\rho:G_n\to\operatorname{Isom}(X)\). Under the compact-quotient convention, it is cocompact if the orbit space \(X/\rho(G_n)\), with the quotient topology, is compact. Another convention requires a compact subset \(K\subseteq X\) such that \(\rho(G_n)K=X\). The supplied record does not specify its convention or hypotheses ensuring agreement between these formulations.

\textbf{Classification question.} Classify the \(\operatorname{CAT}(0)\) spaces admitting a cocompact action of \(\operatorname{Isom}(\mathbb H^n)\), and classify the \(\operatorname{CAT}(-1)\) spaces admitting such an action.

\textbf{Unresolved qualifications.} The isometric-action interpretation above is the natural geometric reading, but the supplied statement does not explicitly specify the action category. It also does not state whether actions must be continuous for the usual topology on \(G_n\), faithful, or proper, or whether spaces must be complete, proper, locally compact, or geodesically complete. None of these additional hypotheses is imposed here. The source does not specify whether the objects to classify are spaces or spaces equipped with chosen actions, nor whether classification is up to isometry, equivariant isometry, scaling, or another equivalence. Thus the displayed question preserves the source's classification request but does not determine a unique fully specified classification problem.
\subsection{Short English statement}
Which \(\operatorname{CAT}(0)\) or \(\operatorname{CAT}(-1)\) spaces admit a cocompact action of the full isometry group of hyperbolic \(n\)-space? The dimension range, action hypotheses, cocompactness convention, and equivalence relation for classification remain unspecified. The supplied record labels the problem open; this draft does not independently verify that status.
\subsection{Sources}
\begin{itemize}
\item Supplied local registry record, permanent UnsolvedMath ID \(30002140\), title ``Cocompact Hyperbolic Actions on CAT Spaces,'' in the Oberwolfach Reports Open Problems collection; exact source problem code \texttt{catalogue entry}. The record gives proposed year \(2012\) and identifies no proposer.
\item Supplied report citation: \emph{Geometric Group Theory, Hyperbolic Dynamics and Symplectic Geometry} (2013), Oberwolfach Reports, pp.\ 2139--2203, DOI \texttt{10.4171/owr/2012/35}. \url{https://doi.org/10.4171/owr/2012/35}. The background describes Oberwolfach Reports as documenting mathematical workshops and their open-problem discussions. The citation year \(2013\) and the record's proposed year \(2012\) are retained as supplied.
\item Supplied literature source: Monod--Py, ``exotic hyperbolic deformations.'' \url{https://arxiv.org/abs/1201.0606}. This is the supplied descriptive label; no fuller bibliographic identification is provided.
\item Supplied literature source: ``Oberwolfach Report 35.'' \url{https://doi.org/10.4171/owr/2012/35}.
\item Literature assessment dated \texttt{2026-08-21}: the supplied curation labels the problem open and reports that exotic CAT-negative deformations and classifications under additional hypotheses are known, while no classification of all CAT(0) or CAT-negative spaces with cocompact hyperbolic isometry action is available. This is supplied dated curation, not an independently verified literature assessment. The background explicitly requests independent verification of time-sensitive status and cited claims.
\item The registry labels the statement \texttt{corrected\_verified} and reports that its mathematical content was checked against the official Oberwolfach report, with notation, context, and question form normalized without changing the source-level claim. This verification description is supplied metadata; no source visit or independent verification is asserted here.
\end{itemize}
\end{document}
